
\chapter{Внутренний реестр обязательств и событийная модель}\label{ch:internal-ledger}

Шлюз отвечает продавцу за сотни миллисекунд, но в~этот момент
ещё не ясно, произойдёт ли подтверждение (capture), когда придёт
расчёт, какой будет комиссия эквайера и~не превратится ли успешная
продажа через месяц в~чарджбэк. Между пограничным слоем шлюза
и операционной сверкой нужен ещё один слой -- внутренний реестр
обязательств, в~котором продукт фиксирует собственные денежные
обязательства и привязывает их к~событиям платёжного
потока~\cite{visa-core-rules,mc-tpr,stripe-manual-capture}.

\section{Почему шлюз не может быть системой учёта}\label{sec:ledger-why}

Шлюз отвечает на вопрос \enquote{что произошло на границе системы}:
пришёл запрос, сработала проверка риска, эквайер вернул одобрение
или отказ, продавцу ушёл webhook. Для интеграции этого достаточно,
для финансового учёта -- нет, и~причин ровно две.

Во-первых, статус шлюза не равен состоянию денег. Авторизация
со статусом \enquote{approved} означает лишь, что эмитент зарезервировал
сумму или одобрил её к~дальнейшей обработке. Подтверждение, клиринг,
расчёт, отмены и корректировки в~правилах карточных схем идут
отдельными стадиями. Пока сумма не перешла в~доступный остаток,
считать её готовой к~выплате
рано~\cite{visa-core-rules,mc-tpr,stripe-manual-capture,stripe-balances}.

Во-вторых, за одним внешним статусом скрывается несколько денежных
последствий. Успешная транзакция в~маркетплейсе порождает отдельные
обязательства перед ТСП, платформой и эквайером. Хранить это одной
строкой \texttt{payment.status = succeeded} -- значит рано или поздно
потерять связь между движением денег и жизненным циклом продукта.

Маркетплейс делает это нагляднее. Покупатель платит одну сумму,
а~из неё возникают минимум три обязательства: выплата продавцу,
комиссия платформы и удержание под логистику при наличии доставки.
Если платформа гарантирует возврат, добавляется четвёртое --
резерв под возвраты. В~шлюзе всё это одна строка
\texttt{payment = 100}. В~реестре -- четыре проводки с~разными
сроками признания и разными правилами отмены.

\begin{figure}[tbp]
  \centering
  \includefiguresvg[width=.8\linewidth]{assets/figures/ch28-ledger/gateway-vs-ledger}
  \caption{Шлюз vs.\ реестр обязательств: один платёж, четыре проводки.}\label{fig:gateway-vs-ledger}
\end{figure}

\definitionnote{Внутренний реестр обязательств}{Внутренняя учётная
  система платёжного продукта, где каждая операция записывается как
  набор денежных проводок. Реестр фиксирует собственные обязательства
  продукта, работая независимо от банковского счёта и~внешней
бухгалтерии.}

Шлюз из главы~\ref{ch:gateway} отвечает за приём и маршрутизацию.
Реестр отвечает за состояние денег внутри продукта: что уже
заработано, что зарезервировано, что можно выплатить продавцу.

\section{Двойная запись: минимум для платёжного продукта}\label{sec:ledger-double-entry}

Без минимума бухгалтерских терминов дальше читать нельзя.
\textbf{Счёт} -- именованный накопитель денежных значений:
\texttt{cardholder\_receivable}, \texttt{merchant\_pending},
\texttt{platform\_fee\_earned}. \textbf{Проводка} (бухгалтерская
запись) -- пара \enquote{дебет счёта~A, кредит счёта~B на одну
и~ту же сумму}: один счёт получает прирост, другой -- уменьшение,
в~зависимости от типа счёта. \textbf{Двойная запись} (double-entry
bookkeeping) фиксирует каждое денежное событие минимум одной
парной проводкой так, что сумма дебетов всегда равна сумме
кредитов. Расхождение сумм означает ошибку в~журнале -- инвариант
встроен в~саму модель.

Двойная запись отличает реестр обязательств от обычного журнала.
Журнал событий говорит \enquote{покупатель оплатил заказ}; реестр
говорит \enquote{из этого следуют такие движения по счетам},
и~движения автоматически балансируются.

В~платёжном продукте двойная запись решает три задачи. Первая --
делает аудит-след самопроверяющимся: \enquote{потерять} деньги
без видимого расхождения невозможно. Вторая -- разделяет
обязательства и фактические остатки (см.\ pending vs.\ available
ниже). Третья -- упрощает сверку с~внешним миром: банковская
выписка устроена по той же двойной записи со стороны банка,
а сопоставление двух согласованных картин надёжнее сравнения
журнала с~плоской цифрой.

\section{Продуктовый реестр обязательств и бухгалтерский учёт банка}\label{sec:ledger-vs-accounting}

Внутренний реестр платёжного продукта и бухгалтерский учёт
кредитной организации, в~чьём контуре он работает, -- два разных
слоя. Их легко перепутать, особенно если продукт встроен в~банк.
Различия существенны и для архитектуры, и~для регуляторики.

\textbf{Продуктовый реестр обязательств} -- внутренний инструмент
платёжного продукта. Он живёт в~базах данных продукта, использует
произвольную нотацию счетов (\texttt{merchant\_pending},
\texttt{platform\_fee\_earned}, \texttt{refund\_payable}), считает
балансы партнёров -- продавцов, плательщиков, платформ -- и~отвечает
на вопросы продукта: сколько мы должны ТСП, сколько уже признано
комиссией, какой резерв удержан под возвраты. Наружу как
бухгалтерский отчёт он не публикуется.

\textbf{Бухгалтерский учёт банка по РСБУ.} Кредитная организация
в~России ведёт учёт операций по российским стандартам бухгалтерского
учёта (РСБУ) на основе Плана счетов кредитных организаций,
утверждённого Положением Банка России 809-П \enquote{О~Плане счетов
бухгалтерского учёта для кредитных организаций и порядке его
применения}~\cite{cbr-809p}.\footnote{Точные реквизиты действующего
Положения и редакция Плана счетов сверяются по сайту Банка России
перед публикацией; 809-П заменило ранее действовавшее Положение
№ 579-П.} План счетов задаёт фиксированную нумерацию: \enquote{30102} --
корреспондентские счета кредитных организаций в~Банке России,
\enquote{40817} -- счета физических лиц, \enquote{40702} -- счета
коммерческих организаций, \enquote{47422} -- \enquote{Обязательства
по прочим операциям}, на котором, в~частности, учитываются суммы
клиентских платежей в~пути. Банк обязан отражать платёжные операции
на этих счетах в~установленном порядке; содержимое подаётся
в~формальной отчётности Банку России и проверяется аудитом.

Соотношение между двумя слоями устроено так. Продуктовый реестр --
детальный взгляд на одну операцию через её жизненный цикл внутри
продукта; бухгалтерский учёт -- общий взгляд через банковские счета
и обязательные группировки. Каждая запись продуктового реестра
в~конечном счёте порождает одну или несколько бухгалтерских проводок
на счетах 809-П; обратное верно не всегда -- бухгалтерия агрегирует
продуктовые записи и~хранит их сжато. Для платёжного продукта в~банке
принципиально не дублировать двойную запись в~обоих слоях, а~явно
зафиксировать соответствие между внутренними счетами продуктового
реестра и~счетами 809-П, чтобы сверка шла по этому соответствию,
а~не пересобирала бухгалтерию заново.

Если платёжный продукт работает не на стороне банка (например,
в~составе небанковской платформы), связь смещается: бухгалтерия
по РСБУ ведётся у~банка-партнёра или нескольких партнёров,
а продуктовый реестр сцеплен с~ней через интерфейсные документы,
реестры и банковские выписки. В~обоих случаях соответствие между
внутренним реестром и~Планом счетов нужно закладывать в~архитектуру
с~самого начала -- иначе получаются два независимых параллельных
учёта, которые потом не сойдутся.

\section{События: авторизация, подтверждение, возврат, спор}\label{sec:ledger-events}

Реестр строится вокруг событий. В~карточном потоке минимально
нужны события авторизации, подтверждения, возврата, чарджбэка
и начисления комиссии. Они образуют причинную цепочку: каждое
событие меняет деньги по-своему и~не должно растворяться в~общем
статусе \enquote{платёж успешен}~\cite{mc-tpr,mc-chargeback-guide,adyen-glossary}.

Реестр атомарно записывает каждое событие в~журнал и~сразу
порождает проводки. Продуктовая система получает две проекции
одной реальности: \textit{что случилось} и \textit{как это
изменило деньги}.

Авторизация на 5\,000~руб. ещё не создаёт выручку ТСП, но фиксирует
ожидаемое обязательство: при успешном подтверждении сумма перейдёт
в~баланс, доступный к~выплате. Подтверждение превращает ожидание
в~финансовый факт. Возврат запускает обратное движение уже после
подтверждения -- платёж, не дошедший до подтверждения, обычно
отменяют, а~не возвращают~\cite{adyen-glossary,stripe-manual-capture}.
Чарджбэк -- самостоятельное событие: его инициирует банк-эмитент,
оно живёт по правилам платёжной системы и~несёт собственные сроки
и~комиссии (см.\ главу~\ref{ch:disputes})~\cite{adyen-glossary,mc-chargeback-guide}.

\practicenote{%
  Кодируйте возврат и~чарджбэк как отдельные типы операций.
  Для клиента результат похож, но для продукта это разные
  процессы: возврат инициирует продавец, а~чарджбэк приходит извне и
  живёт в~цикле оспаривания.
}

Событие отвечает за причинность, проводка -- за денежный смысл.
Попытка восстановить одно из другого задним числом быстро лишает
продукт объяснимости.

Переходы между событиями образуют конечный автомат. Авторизация
переходит в~подтверждение или в void, но не наоборот. Подтверждение
переходит в~рассчитанную; после расчёта возможен возврат или чарджбэк.
Void допустим только до подтверждения. Возврат после расчёта порождает
обратное движение средств; частичный возврат уменьшает обязательство,
не отменяя его. Чарджбэк приходит извне и~может наступить после
расчёта, открывая отдельную цепочку -- повторное представление,
pre-arbitration, arbitration~\cite{visa-core-rules,mc-chargeback-guide}.

Каждый переход порождает проводку. Auth создаёт pending; подтверждение
переводит pending в earned; void обнуляет pending; расчёт переводит
earned в available; возврат дебетует available и кредитует
payable-to-cardholder; чарджбэк дебетует available или reserve
и открывает отдельный учётный цикл спора. Без явно описанного автомата
команда кодирует переходы в~разрозненных условиях, и~через год никто
уже не помнит, допустим ли возврат после частичного чарджбэка.

Продавец продаёт товар за 5\,000~руб. через маркетплейс с~комиссией
платформы 10\,\%.

\textbf{Auth (день~0).} Эмитент одобряет авторизацию. Реестр
создаёт две проводки: дебет \texttt{cardholder\_receivable}
5\,000~руб., кредит \texttt{merchant\_pending} 5\,000~руб.
Деньги ещё не у~платформы, но обязательство зафиксировано.

\textbf{Подтверждение (день~0, через 30~мин).} Продавец подтверждает
отгрузку. Реестр переводит pending в earned: дебет
\texttt{merchant\_pending} 5\,000~руб., кредит
\texttt{merchant\_earned} 4\,500~руб. и~кредит
\texttt{platform\_fee\_earned} 500~руб. С~этого момента
платформа признала свою комиссию.

\textbf{Расчёт (день~2).} Эквайер перечисляет средства.
Реестр переводит earned в available: дебет
\texttt{merchant\_earned} 4\,500~руб., кредит
\texttt{merchant\_available} 4\,500~руб. Продавец может
запросить выплату.

\textbf{Частичный возврат (день~5).} Покупатель возвращает часть
товара на 2\,000~руб. Реестр: дебет \texttt{merchant\_available}
2\,000~руб., кредит \texttt{refund\_payable} 2\,000~руб.
Комиссия платформы не пересчитывается (политика маркетплейса);
если бы пересчитывалась, отдельная проводка уменьшила бы
\texttt{platform\_fee\_earned} на 200~руб.

Каждый переход оставляет аудиторский след: сумма дебетов
всегда равна сумме кредитов, а~баланс любого счёта
восстанавливается простым сканом журнала проводок.

\begin{figure}[tbp]
  \centering
  \includefiguresvg[width=.8\linewidth]{assets/figures/ch28-ledger/event-state-machine}
  \caption{Конечный автомат событий: от авторизации до чарджбэка.}\label{fig:event-state-machine}
\end{figure}

\section{Балансы, комиссии и~правила перехода}\label{sec:ledger-balances}

Платёжный продукт оперирует двумя состояниями одного баланса
и отдельным механизмом резерва.

\textbf{Pending} и \textbf{available} -- два состояния одного и~того
же основного баланса. Pending -- сумма, недоступная к~выплате:
она находится между авторизацией и подтверждением, ждёт клиринга
или удерживается внутри окна выплат. Available -- та часть, которой
платформа готова распоряжаться: перечислить продавцу, зачесть
в~счёт комиссии или использовать для внутреннего неттинга. Деньги
последовательно переходят из pending в available по мере прохождения
этапов жизненного цикла операции~\cite{stripe-balances,adyen-balances}.

\textbf{Reserve / holdback} -- отдельный балансовый механизм,
не очередное состояние того же баланса. Reserve хранит сумму,
изъятую из available по правилу риска (refund reserve, minimum
balance для покрытия будущих возвратов, dispute и fees или
отдельный holdback под ожидаемые списания), и подчиняется собственным
правилам пополнения, разморозки и закрытия~\cite{stripe-balances,stripe-minimum-balances,adyen-balances}.

Одна внешняя транзакция обычно порождает несколько проводок.
Подтверждение дебетует \enquote{acquirer receivable} и кредитует
\enquote{merchant pending}. После расчёта реестр переводит сумму
из \enquote{merchant pending} в \enquote{merchant available},
одновременно начисляя комиссию платформы на отдельный счёт.

Без явного моделирования этих переходов продукт начинает путать
\enquote{деньги есть} с \enquote{деньги когда-нибудь придут}.
Отсюда и расхождения между панелью продавца и банковской выпиской.

Помимо основных сумм платежей, реестр фиксирует все производные
денежные эффекты: когда платформа вправе признать свою комиссию
и~когда обязана отложить часть денег под риск. Без этого продукт
принимает финансовые решения по косвенным сигналам из интеграций.

Комиссия платформы редко возникает в~один момент с~авторизацией.
Обычно её признают при подтверждении или расчёте, когда уже ясно,
что операция состоялась. Эквайерские сборы, сетевые комиссии, валютная
наценка и~комиссии за чарджбэк приходят со своими сроками, отличными
от срока основной суммы. Реестр обязан поддерживать логику начислений:
сначала признать ожидаемую комиссию, затем скорректировать её
по фактическому расчётному файлу~\cite{visa-core-rules,mc-tpr,stripe-balances}.

То же касается резервов. Резерв возвратов, minimum balance или любой
другой удерживаемый резерв -- это конкретное правило перевода части
available в~отдельный баланс с~собственным сроком разморозки.

\begin{figure}[tbp]
  \centering
  \includefiguresvg[width=.7\linewidth]{assets/figures/ch28-ledger/balance-states}
  \caption{Три состояния баланса: pending, available, reserve.}\label{fig:balance-states}
\end{figure}

Не выраженное в~реестре правило начинают повторять вручную
в~процессе выплат и~в~бэк-офисе, и~система теряет единый источник
истины~\cite{stripe-minimum-balances,adyen-balances}.

График выплат тоже лучше моделировать в~реестре, а~не в~коде
экспортёра банковских файлов. Ручной режим, ежедневные выплаты,
еженедельные и ежемесячные окна перечислений, задержка T+X --
всё это правила перевода между внутренними балансами~\cite{stripe-payouts}.

Если комиссии, резервы и~график выплат живут вне реестра, реестр
перестаёт быть целостной картиной обязательств.

\section{Идемпотентность и передача в~сверку}\label{sec:ledger-idempotency}

Шлюз и~реестр живут в~разных временных режимах. Шлюз обязан быстро
вернуть продавцу ответ по API. Реестр обязан отразить событие ровно
один раз, даже если внешний мир несколько раз повторил один и~тот же
сигнал~\cite{stripe-idempotency,adyen-api-idempotency}.

Границ идемпотентности две. Первая -- на входе в~шлюз: продавец
повторяет запрос с~тем же \texttt{idempotency-key}, и~шлюз не должен
создать вторую авторизацию. Вторая -- между шлюзом и~реестром:
webhook от эквайера, повторный обратный вызов после тайм-аута или
заново проигранный пакет не должны породить дубль
проводок~\cite{stripe-idempotency,adyen-api-idempotency,stripe-webhooks}.

Реализация устроена последовательно. Сначала шлюз присваивает
операции собственный стабильный идентификатор. Затем привязывает
к~нему все внешние сообщения через таблицу соответствия между STAN,
RRN, идентификатором процессинга и внутренним
ID~\cite{fis-iso8583-guide-2023}. Реестр принимает только
нормализованные внутренние события с~уникальным \path{event_id};
повторная доставка такого события возвращает прежний результат,
не порождая новых проводок~\cite{stripe-webhooks}.

Журнал событий и~денежные проводки разводятся на два шага. Сначала
система фиксирует факт получения сигнала от внешнего мира, затем
атомарно применяет его к~денежным балансам. Сбой между шагами
не страшен -- операцию можно безопасно доиграть.

Сверка из главы~\ref{ch:reconciliation} сопоставляет внутреннее
состояние продукта с~внешними файлами и банковскими движениями.
Для этого продукт обязан сначала зафиксировать собственное состояние.
Шлюз записывает протокольный факт авторизации, реестр превращает его
в~денежное состояние, а~сверка сопоставляет результат с TC33, IPM
и банковской выпиской~\cite{visa-core-rules,mc-tpr,mc-switching-services}.

Без такого реестра сверка пытается напрямую сопоставить внешние
файлы со статусами шлюза, и~команда получает типичный хаос:
\enquote{вроде операция одобрена, но непонятно, это уже доход,
ожидаемая выплата или просто авторизационная блокировка}. Реестр
снимает эту неоднозначность ещё до начала сверки.

\subsection*{Граничные случаи: orphan auth и~гонка возврата с~чарджбэком}

Реестры, спроектированные только под happy path, регулярно ломают
два сценария.

\textbf{Orphan auth.} Авторизация одобрена, реестр создал
\texttt{merchant\_pending}, но подтверждение не пришло: продавец
отменил заказ, потерял webhook или просто забыл. Hold на карте
держателя снимается автоматически по истечении срока, который
правила платёжной системы устанавливают для категории
операции.\footnote{По редакции Visa Authorization Framework
от 13~апреля 2024~года типичные пределы: 5~дней для card-present,
10~дней для card-not-present и~до 30~дней для lodging, cruise lines
и vehicle rental при наличии Estimated Authorization Indicator.
Mastercard и НСПК устанавливают собственные пределы, требующие
отдельной сверки.}~\cite{visa-core-rules,mc-rules,nspk-mir-rules},
а~в~реестре запись \texttt{merchant\_pending} остаётся вечно
без механизма экспирации. Решение -- фоновый процесс, сканирующий
проводки в~ожидании (\textit{pending}) старше порога (максимальный срок hold по категории
плюс буфер) и автоматически порождающий событие
\texttt{auth\_expired}: дебет \texttt{merchant\_pending}, кредит
\texttt{cardholder\_receivable} обнуляют обязательство. Без этого
баланс pending растёт, искажая отчётность и блокируя сверку.

\textbf{Гонка возврата и чарджбэка.} Продавец инициирует возврат
на 5\,000~руб. в~тот же день, когда эмитент открывает чарджбэк
на ту же сумму. Оба события дебетуют \texttt{merchant\_available}.
Обработка обоих списывает у~продавца 10\,000~руб. вместо 5\,000.
Защита -- блокировка на уровне транзакции: перед записью возврата
реестр проверяет, нет ли открытого спора; перед записью чарджбэка --
нет ли уже проведённого возврата. Если возврат уже исполнен, чарджбэк
переходит в~статус \texttt{already\_refunded}, а платформа отправляет
повторное представление с~доказательством
возврата~\cite{mc-chargeback-guide}. Если первым пришёл чарджбэк,
возврат блокируется, и~оператор разрешает конфликт вручную.
Итоговый дебет с~ТСП в~обоих случаях не превышает исходную сумму
транзакции.

\practicenote{%
  Операционной команде нужна таблица соответствия внешних
  идентификаторов и внутренних счетов: авторизация, процессинг, заказ,
  выплата, спор. Без неё даже хороший реестр трудно сверить с~реальным миром.

}
